% Part: computability
% Chapter: computability-theory
% Section: complement-ce

\documentclass[../../../include/open-logic-section]{subfiles}

\begin{document}

\olfileid{cmp}{thy}{cmp}
\olsection{可计算枚举集对补不封闭}

假设 $A$ 是可计算枚举的，那么 $A$ 的补集 $\Complement{A} = \Nat \setminus A$ 是否也总是可计算枚举的？下面的定理和推论表明答案是否定的。

\begin{thm}
\ollabel{thm:ce-comp}
设 $A$ 是任意自然数集合。则 $A$ 是可计算的，当且仅当 $A$ 和 $\Complement{A}$ 都是可计算枚举的。
\end{thm}

\begin{proof}
前向方向是显然的：如果 $A$ 是可计算的，那么 $\Complement{A}$ 也是可计算的（$\Char{A} = 1 \tsub
\Char{\Complement{A}}$），因此两者都是可计算枚举的。

反之，假设 $A$ 和 $\Complement{A}$ 都是可计算枚举的。设 $A$ 为 $\cfind{d}$ 的定义域，$\Complement{A}$ 为 $\cfind{e}$ 的定义域。如下定义 $h$：
\[
h(x) = \umin{s}{(T(d,x,s) \lor T(e,x,s))}.
\]
换言之，对于输入 $x$，$h$ 会搜索 $\cfind{d}$ 或 $\cfind{e}$ 的停机计算。此时，如果 $x \in A$，前一种搜索将会成功；如果 $x \in \Complement{A}$，后一种搜索将会成功。因此，$h$ 是全可计算函数。由此可知，对每个 $x$，$x \in A$ 当且仅当 $T(e, x, h(x))$，即有定义的是 $\cfind{e}$。由于 $T(e, x, h(x))$ 是可计算关系，因此 $A$ 是可计算的。
\end{proof}

\begin{explain}
用非形式的计算术语更容易理解这里的直观：要判定 $A$，只需以 $x$ 为输入，搜索 $\cfind{e}$ 和 $\cfind{f}$ 的停机计算。两者之中必然有一个会停机；如果是 $\cfind{e}$ 停机，那么 $x \in A$，否则 $x \in \Complement{A}$。
\end{explain}

\begin{cor}
\ollabel{cor:comp-k}
$\Complement{K_0}$ 不是可计算枚举的。
\end{cor}

\begin{proof}
已知 $K_0$ 是可计算枚举的但不可计算。如果 $\Complement{K_0}$ 也是可计算枚举的，那么根据 \olref{thm:ce-comp}，$K_0$ 就是可计算的，这与 \olref[ncp]{thm:K-0} 矛盾。
\end{proof}

\end{document}